%% LyX 2.3.7 created this file.  For more info, see http://www.lyx.org/.
%% Do not edit unless you really know what you are doing.
\documentclass[english]{beamer}
\usepackage{mathptmx}
\renewcommand{\sfdefault}{lmss}
\usepackage[T1]{fontenc}
\usepackage[latin9]{inputenc}
\usepackage{amsbsy}
\usepackage{amssymb}
\usepackage{stackrel}
\usepackage{graphicx}

\makeatletter
%%%%%%%%%%%%%%%%%%%%%%%%%%%%%% Textclass specific LaTeX commands.
% this default might be overridden by plain title style
\newcommand\makebeamertitle{\frame{\maketitle}}%
% (ERT) argument for the TOC
\AtBeginDocument{%
  \let\origtableofcontents=\tableofcontents
  \def\tableofcontents{\@ifnextchar[{\origtableofcontents}{\gobbletableofcontents}}
  \def\gobbletableofcontents#1{\origtableofcontents}
}

%%%%%%%%%%%%%%%%%%%%%%%%%%%%%% User specified LaTeX commands.
\usetheme{Boadilla}
\usecolortheme{default}
\setbeamercovered{transparent}
\usepackage{indentfirst}
\usepackage{amsmath,amsfonts,amsthm,bm}

\makeatother

\usepackage{babel}
\begin{document}
\title{\textbf{Derivation of SRTP-DKH Hamiltonian and 2-component electron correlation}}
\author{Liu Kunyu}
\date{04/05/24}
\makebeamertitle
\begin{frame}{Contents}
\begin{description}
\item [{1}] \textbf{Derivation of SRTP-fpFW Hamiltonian}
\begin{description}
\item [{1.1}] \textbf{Operatorization of SRTP Hamiltonian}
\item [{1.2}] \textbf{Back to SRTP-fpFW Hamiltonian}
\end{description}
\item [{\textcolor{lightgray}{2}}] \textcolor{lightgray}{DKH transformation
of }\textbf{SRTP-fpFW}\textcolor{lightgray}{{} Hamiltonian}
\begin{description}
\item [{\textcolor{lightgray}{2.1}}] \textbf{SRTP-DKH2-1e Hamiltonian}
\item [{\textcolor{lightgray}{2.2}}] \textcolor{lightgray}{Operatorization
of SRTP Hamiltonian}
\end{description}
\item [{\textcolor{lightgray}{3}}] \textcolor{lightgray}{Modified DKH Hamiltonian}
\begin{description}
\item [{\textcolor{lightgray}{3.1}}] \textcolor{lightgray}{Operatorization
of SRTP Hamiltonian}
\item [{\textcolor{lightgray}{3.2}}] \textcolor{lightgray}{Operatorization
of SRTP Hamiltonian}
\end{description}
\begin{block}{Agreements}
\begin{itemize}
\item {\small{}All derication is based on a.u., the charge and mass of electron
is equal to 1, the magnitude of spin angular momentum is equal to
1/2. }{\small\par}
\item {\small{}Classical vector is represented as $\overrightarrow{\mathrm{a}}$;
vector operator is represented as $\mathbf{\hat{\mathbf{a}}}$.}{\small\par}
\end{itemize}
\end{block}
\end{description}
\end{frame}
\pagebreak{}
\begin{frame}{1.1 Operatorization of SRTP Hamiltonian}
Second Relativized Thomas Precession (SRTP) Hamiltonian of single
spin-1/2 electron is

{\footnotesize{}
\begin{equation}
H_{\mathrm{SRTP}}=\frac{\gamma}{2\beta^{2}}\overrightarrow{S}_{\gamma}\cdot(\dot{\overrightarrow{\beta}}\times\overrightarrow{\beta})\frac{\frac{4}{3}\gamma^{2}(\overrightarrow{\beta}\cdot\overrightarrow{S})^{2}}{(1+\frac{4}{3}\gamma^{2}(\overrightarrow{\beta}\cdot\overrightarrow{S})^{2})^{\frac{3}{2}}}
\end{equation}
}{\footnotesize\par}
this expression is fully classic. {\footnotesize{}$\gamma,\overrightarrow{\beta},\overrightarrow{S},$}\textrm{\footnotesize{}$\dot{\overrightarrow{\beta}}$}
is the Lorentz factor, velocity, spin angular momentum and coordinate
acceleration of electron, respectively. \textrm{\footnotesize{}$S_{\gamma,i}$}
is a vector related to spin and motion of electron:

{\footnotesize{}
\begin{equation}
S_{\gamma,\mathrm{i}}=\frac{1}{\sqrt{1-\beta_{\mathrm{i}}^{2}}}S_{\mathrm{i}}
\end{equation}
}For the next step, a kinetic transformation is performed:{\footnotesize{}
\begin{align*}
\dot{\overrightarrow{\beta}}\times\overrightarrow{\beta} & =\frac{1}{\gamma c}\left(-\overrightarrow{\nabla}(V)+\overrightarrow{\beta}(\overrightarrow{\beta}\cdot\overrightarrow{\nabla}(V))\right)\times\overrightarrow{\beta}\\
 & =-\frac{1}{\gamma^{2}c^{2}}(\overrightarrow{\nabla}(V)\times\overrightarrow{p})
\end{align*}
}{\footnotesize\par}

\end{frame}
\pagebreak{}
\begin{frame}{1.1 Operatorization of SRTP Hamiltonian}
For {\footnotesize{}$\overrightarrow{S}_{\gamma}$}, we consider $\widehat{\beta}_{\mathrm{i}}$in
its contraction factor as the group velocity without Zitterbewegung

{\footnotesize{}
\begin{equation}
\widehat{S}_{\gamma,\mathrm{i}}=\frac{1}{2}\frac{1}{\sqrt{1-\widehat{\beta}_{\mathrm{i}}^{2}}}\widehat{\sigma}_{\mathrm{i}}=\frac{1}{2}\frac{1}{\sqrt{1-\frac{c^{2}}{\epsilon^{2}}\widehat{p}_{\mathrm{i}}^{2}}}\widehat{\sigma}_{\mathrm{i}}
\end{equation}
}where \textrm{$\epsilon$ }is the energy of electron, treatment of
\textrm{$\epsilon$} will be discussed in next section. Then SRTP
Hamiltonian can be further operatorized as:

{\footnotesize{}
\begin{equation}
\widehat{H}_{\mathrm{SRTP}}=-\frac{1}{2}\left\{ \frac{\epsilon(\frac{\widehat{\boldsymbol{p}}}{c}\cdot\widehat{\boldsymbol{\sigma}})^{2}}{2c^{4}(\frac{\widehat{\boldsymbol{p}}}{c})^{2}\left(1+(\frac{\widehat{\boldsymbol{p}}}{c}\cdot\widehat{\boldsymbol{\sigma}})^{2}\right)^{\frac{3}{2}}},\widehat{\boldsymbol{S}}_{\gamma}\cdot(\boldsymbol{\nabla}(\widehat{V})\times\widehat{\boldsymbol{p}})\right\} 
\end{equation}
}we ignore the contribution of magnetic field to SRTP term, it will
be partially compensated by the regularization of {\footnotesize{}$\widehat{\boldsymbol{p}}$}
and {\footnotesize{}$\epsilon$}. This operatorization refers to the
classical Thomas precession Hamiltonian obtained from the fpFW transformation
of Dirac equation. (J. Math. Phys., 2003, 44(7): 2952--2966) Furthermore,
fpFW is able to give an accurate result of classical Thomas precession
Hamiltonian, and more suitable for bound states {\footnotesize{}$(\lambda_{\mathrm{B}}\sim l_{\mathrm{c}})$}
compared to iterative FW transformation. (J. Math. Phys., 2013, 176(2):
987-999)

This form of the SRTP Hamiltonian doesn't have a corresponding retarded
term, since it only contains electric interaction, which is instantaneous
interaction in Coulomb gauge. 

\end{frame}
\pagebreak{}
\begin{frame}{1.1 Operatorization of SRTP Hamiltonian }

Additionally, mathematical form of {\footnotesize{}$\widehat{H}_{\mathrm{SRTP}}$}
allows for no retarded term, similar to magnetic interaction spin-other-orbital
coupling term in fpFW Hamiltonian.{\footnotesize{}
\[
\widehat{H}_{\mathrm{SOC3-ret}}=\underset{\mu<\nu}{\sum}-\frac{2}{c^{2}r_{\mu\nu}^{3}}\frac{\left(\widehat{\boldsymbol{S}}_{\upsilon}\cdot\boldsymbol{r}_{\mu\nu}\right)\left(\boldsymbol{r}_{\mu\nu}\times\widehat{\boldsymbol{p}}_{\mu}\right)\cdot\boldsymbol{r}_{\mu\nu}}{r_{\mu\nu}^{2}}=0
\]
}{\footnotesize\par}

By decomposing the potential energy, the SRTP Hamiltonian are distinguished
into one-electron and two-electron terms as{\footnotesize{}
\begin{equation}
\widehat{H}_{\mathrm{SRTP1}}=-\frac{1}{2}\underset{\mathrm{nuc}}{\sum}\left\{ \frac{\epsilon\widehat{h}^{2}}{2c^{4}\left(1+(\frac{\widehat{\boldsymbol{p}}}{c}\cdot\widehat{\boldsymbol{\sigma}})^{2}\right)^{\frac{3}{2}}},\widehat{\boldsymbol{S}}_{\gamma}\cdot(\boldsymbol{\nabla}(\widehat{V})\times\widehat{\boldsymbol{p}})\right\} 
\end{equation}
\begin{equation}
\widehat{H}_{\mathrm{SRTP2}}=\frac{1}{2}\underset{\mathrm{ele}}{\sum}\left\{ \frac{\epsilon\widehat{h}^{2}}{2c^{4}\left(1+(\frac{\widehat{\boldsymbol{p}}}{c}\cdot\widehat{\boldsymbol{\sigma}})^{2}\right)^{\frac{3}{2}}},\widehat{\boldsymbol{S}}_{\gamma}\cdot(\boldsymbol{\nabla}(\widehat{V})\times\widehat{\boldsymbol{p}})\right\} 
\end{equation}
}\textrm{\footnotesize{}$\widehat{h}$}\textrm{ }is the helicity operator,
which eigenvalue is {\footnotesize{}$\pm1$}. Treatment of \textrm{\footnotesize{}$\widehat{h}^{2}$}
will be discussed in next section.

\end{frame}
\pagebreak{}
\begin{frame}{1.1 Operatorization of SRTP Hamiltonian }

Since the denominator of the coefficient contains \textrm{\footnotesize{}$(\frac{\widehat{\boldsymbol{p}}}{c})^{2}$},
although it converge at \textrm{\footnotesize{}$\frac{\overrightarrow{p}}{c}\rightarrow0$},
it's still impossible to perform a power series expansion at \textrm{\footnotesize{}$\frac{\overrightarrow{p}}{c}=0$}.
Fortunately, the helicity operator has an ideal property.
\begin{theorem}
The square of helicity is equal to 1 in any case.{\footnotesize{}
\begin{align*}
\widehat{h}^{2} & =\frac{(\widehat{\boldsymbol{p}}\cdot\widehat{\boldsymbol{S}})^{2}}{p^{2}S^{2}}\\
 & =\stackrel[\mathrm{i}]{3}{\sum}\frac{-\partial_{\mathrm{i}}^{2}\sigma_{\mathrm{i}}^{2}-\stackrel[\mathrm{j<i}]{3}{\sum}\left\{ \partial_{\mathrm{i}}\sigma_{\mathrm{j}},\partial_{\mathrm{j}}\sigma_{\mathrm{i}}\right\} }{-\partial_{\mathrm{i}}^{2}}\\
 & =\stackrel[\mathrm{i}]{3}{\sum}\frac{-\partial_{\mathrm{i}}^{2}-\stackrel[\mathrm{j<i}]{3}{\sum}\left\{ \sigma_{\mathrm{j}},\sigma_{\mathrm{i}}\right\} \partial_{\mathrm{i}}\partial_{\mathrm{j}}}{-\partial_{\mathrm{i}}^{2}}\\
 & =1
\end{align*}
}{\footnotesize\par}
\end{theorem}

This is because the helicity and spin components of each state have
only two observations: 1 and -1, although the expectations are complex.
\end{frame}
\pagebreak{}
\begin{frame}{1.1 Operatorization of SRTP Hamiltonian}
\begin{theorem}
The square of the spin component of each spin state measured from
each direction is 1.

Define \textrm{\footnotesize{}$\left|1\right\rangle $} and \textrm{\footnotesize{}$\left|-1\right\rangle $}
to be the eigenstates of arbitrarily oriented spin component {\footnotesize{}$\widehat{\sigma}_{\mathrm{i}}$},
and \textrm{\footnotesize{}$\left|\chi\right\rangle $} to be arbitrarily
normalized spin state, \textrm{\footnotesize{}$\left|\chi\right\rangle =c_{1}\left|1\right\rangle +c_{2}\left|-1\right\rangle $},
we have{\footnotesize{}
\begin{align*}
\left\langle \chi\right|\widehat{\sigma}_{\mathrm{i}}\left|\chi\right\rangle  & =c_{1}\widehat{\sigma}_{\mathrm{i}}\left|1\right\rangle +c_{2}\widehat{\sigma}_{\mathrm{i}}\left|-1\right\rangle \\
 & =\left(c_{1}\left\langle 1\right|+c_{2}\left\langle -1\right|\right)\left(c_{1}\left|1\right\rangle -c_{2}\left|-1\right\rangle \right)\\
 & =c_{1}^{2}-c_{2}^{2}\\
\left\langle \chi\right|\widehat{\sigma}_{\mathrm{i}}^{2}\left|\chi\right\rangle  & =c_{1}\widehat{\sigma}_{\mathrm{i}}^{2}\left|1\right\rangle +c_{2}\widehat{\sigma}_{\mathrm{i}}^{2}\left|-1\right\rangle \\
 & =\left(c_{1}\left\langle 1\right|+c_{2}\left\langle -1\right|\right)\left(c_{1}\left|1\right\rangle +c_{2}\left|-1\right\rangle \right)\\
 & =c_{1}^{2}+c_{2}^{2}\\
 & =1
\end{align*}
}{\footnotesize\par}

\pagebreak{}
\end{theorem}

Bring above results into Eq. (5), Eq. (6) and simplify:{\footnotesize{}
\begin{align}
\widehat{H}_{\mathrm{SRTP1}} & =-\frac{1}{2}\underset{\mathrm{nuc}}{\sum}\left\{ \frac{c^{2}}{4\epsilon^{2}},\mathop{\stackrel[\mathrm{i=1}]{3}{\sum}}\left(\epsilon_{\mathrm{i}}\widehat{\sigma}_{\mathrm{i}}\frac{\mathrm{Z}}{r^{3}}(\boldsymbol{r}\times\widehat{\boldsymbol{p}})_{\mathrm{i}}\right)\right\} \nonumber \\
 & =-\frac{1}{4}\underset{\mathrm{nuc}}{\sum}\left\{ \frac{c^{2}}{4\epsilon^{2}},\widehat{\mathfrak{\boldsymbol{s}}}\cdot\frac{\mathrm{Z}}{r^{3}}(\boldsymbol{r}\times\widehat{\boldsymbol{\pi}})-\widehat{\mathfrak{\boldsymbol{s}}}\cdot\frac{\mathrm{Z}}{r^{3}}(\widehat{\boldsymbol{\pi}}\times\boldsymbol{r})\right\} \\
\widehat{H}_{\mathrm{SRTP2}} & =\frac{1}{2}\underset{\mathrm{ele}}{\sum}\left\{ \frac{c^{2}}{4\epsilon^{2}},\mathop{\stackrel[\mathrm{i=1}]{3}{\sum}}\left(\epsilon_{\mathrm{i}}\widehat{\sigma}_{\mathrm{i}}\frac{\mathrm{Z}}{r^{3}}(\boldsymbol{r}\times\widehat{\boldsymbol{p}})_{\mathrm{i}}\right)\right\} \nonumber \\
 & =\frac{1}{4}\underset{\mathrm{ele}}{\sum}\left\{ \frac{c^{2}}{4\epsilon^{2}},\widehat{\mathfrak{\boldsymbol{s}}}\cdot\frac{1}{r^{3}}(\boldsymbol{r}\times\widehat{\boldsymbol{\pi}})-\widehat{\mathfrak{\boldsymbol{s}}}\cdot\frac{1}{r^{3}}(\widehat{\boldsymbol{\pi}}\times\boldsymbol{r})\right\} 
\end{align}
}which {\footnotesize{}$\widehat{\mathfrak{s}}_{\mathrm{i}}=\epsilon_{\mathrm{i}}\widehat{\sigma}_{\mathrm{i}}$}
is Second Relativized Spin (SRS). Consider the commutate relation
{\footnotesize{}$\left[\mathrm{f}\left(\widehat{p}^{2},\widehat{p}_{\mathrm{i}}\right),\widehat{L}_{\mathrm{i}}\right]=0$},
SRS-type SOC is still a self-adjoint operator.{\footnotesize{}
\[
\left(\widehat{\mathfrak{s}}\cdot\widehat{\boldsymbol{L}}\right)^{\dagger}=\widehat{\boldsymbol{L}}\cdot\widehat{\mathfrak{s}}=\widehat{\mathfrak{s}}\cdot\widehat{\boldsymbol{L}}
\]
}This form of SRTP Hamiltonian is the bridge connects modified fpFW
Hamiltonian and Eq. (1). For further steps, we have to admit Eq. (8)
applies to the whole odd term contains canonical momentum, which means
\textrm{\footnotesize{}$\widehat{\boldsymbol{p}}_{\mathrm{i}}$} in
Eq. (8) will be replaced by\textrm{ }\textrm{\footnotesize{}$\hat{\mathbf{\sigma}}_{\mathrm{i}}\widehat{\boldsymbol{\pi}}_{\mathrm{i}}$}
(ignore the coefficient of $\hat{\mathbf{\sigma}}_{\mathrm{i}}$)\textrm{
}and introduce Darwin term{\footnotesize{}
\begin{align}
\widehat{H}_{\mathrm{SRTP1}} & =-\frac{1}{4}\underset{\mathrm{nuc}}{\sum}\left\{ \frac{c^{2}}{4\epsilon^{2}},\widehat{\mathfrak{\boldsymbol{s}}}\cdot\frac{\mathrm{Z}}{r^{3}}(\boldsymbol{r}\times\widehat{\boldsymbol{\pi}})-\widehat{\mathfrak{\boldsymbol{s}}}\cdot\frac{\mathrm{Z}}{r^{3}}(\widehat{\boldsymbol{\pi}}\times\boldsymbol{r})-\overrightarrow{\boldsymbol{\nabla}}\cdot\overrightarrow{\boldsymbol{E}}\right\} \\
\widehat{H}_{\mathrm{SRTP2}} & =\frac{1}{4}\underset{\mathrm{ele}}{\sum}\left\{ \frac{c^{2}}{4\epsilon^{2}},\widehat{\mathfrak{\boldsymbol{s}}}\cdot\frac{1}{r^{3}}(\boldsymbol{r}\times\widehat{\boldsymbol{\pi}})-\widehat{\mathfrak{\boldsymbol{s}}}\cdot\frac{1}{r^{3}}(\widehat{\boldsymbol{\pi}}\times\boldsymbol{r})-\overrightarrow{\boldsymbol{\nabla}}\cdot\overrightarrow{\boldsymbol{E}}\right\} 
\end{align}
}magnetic SOC term is related to kinetic term, follow the form $\left\{ \frac{1}{\epsilon},\widehat{\boldsymbol{\sigma}}\cdot\overrightarrow{\boldsymbol{H}}\right\} $,
it's two-electron term with $\overrightarrow{\boldsymbol{H}}$ contains
retarded magnetic interaction between electrons in absence of external
magnetic field.

Additionally, we have to change the coefficient in Eq. (9), (10),
because the odd term will be zero and the energy-momentum relationship
will be correct if and only if we use the transformation $U^{\pm}=\frac{\epsilon+c^{2}\pm\beta\mathcal{O}}{\sqrt{2\epsilon\left(\epsilon+c^{2}\right)}}$,
it gives the coefficient as:

{\footnotesize{}
\begin{align}
\widehat{H}_{\mathrm{SRTP1}} & =-\frac{1}{8}\underset{\mathrm{nuc}}{\sum}\left\{ \frac{c^{2}}{\epsilon(\epsilon+c^{2})},\widehat{\mathfrak{\boldsymbol{s}}}\cdot\frac{\mathrm{Z}}{r^{3}}(\boldsymbol{r}\times\widehat{\boldsymbol{\pi}})-\widehat{\mathfrak{\boldsymbol{s}}}\cdot\frac{\mathrm{Z}}{r^{3}}(\widehat{\boldsymbol{\pi}}\times\boldsymbol{r})-\overrightarrow{\boldsymbol{\nabla}}\cdot\overrightarrow{\boldsymbol{E}}\right\} \\
\widehat{H}_{\mathrm{SRTP2}} & =\frac{1}{8}\underset{\mathrm{ele}}{\sum}\left\{ \frac{c^{2}}{\epsilon(\epsilon+c^{2})},\widehat{\mathfrak{\boldsymbol{s}}}\cdot\frac{1}{r^{3}}(\boldsymbol{r}\times\widehat{\boldsymbol{\pi}})-\widehat{\mathfrak{\boldsymbol{s}}}\cdot\frac{1}{r^{3}}(\widehat{\boldsymbol{\pi}}\times\boldsymbol{r})-\overrightarrow{\boldsymbol{\nabla}}\cdot\overrightarrow{\boldsymbol{E}}\right\} 
\end{align}
}{\footnotesize\par}

this result can be seen as the conbination of the coefficient given
by $\frac{\partial U}{\partial\theta}$ and the SRS given by $\frac{\partial U}{\partial\varphi}$.

It's worth to mention that the treatment of {\footnotesize{}$\epsilon$},
if use the approximation {\footnotesize{}$\epsilon=\sqrt{c^{4}+\widehat{p}^{2}c^{2}}$},
Eq. (9) and Eq. (10) contains only electronic and magnetic SOC terms;
else if {\footnotesize{}$\epsilon=\sqrt{c^{4}+(\hat{\mathbf{\sigma}}\cdot\widehat{\boldsymbol{\pi}})^{2}c^{2}}$,
}Eq. (9) and Eq. (10) contains also electronic-magnetic crossover
SOC terms. 
\end{frame}
\pagebreak{}
\begin{frame}{1.2 \textbf{Back to SRTP-fpFW Hamiltonian}}
\begin{theorem}
{\footnotesize{}$(\widehat{\boldsymbol{\mathfrak{a}}}\cdot\widehat{\boldsymbol{p}})(\widehat{\boldsymbol{\mathfrak{a}}}\cdot\widehat{\boldsymbol{q}})=\mathrm{i\widehat{\mathfrak{\boldsymbol{s}}}\cdot(\mathit{\widehat{\boldsymbol{q}}}\times\widehat{\boldsymbol{\mathit{p}}})}+\mathrm{spin.free.term}$
}is ture if and only if{\footnotesize{} $\widehat{\mathfrak{a}}_{\mathrm{i}}=\sqrt{\frac{\epsilon_{\mathrm{j}}\epsilon_{\mathrm{k}}}{\epsilon_{\mathrm{i}}}}\widehat{\alpha}_{\mathrm{i}}$.}{\footnotesize\par}

{\footnotesize{}
\begin{align*}
 & (\widehat{\boldsymbol{\mathfrak{a}}}\cdot\widehat{\boldsymbol{p}})(\widehat{\boldsymbol{\mathfrak{a}}}\cdot\widehat{\boldsymbol{q}})\\
= & \widehat{\mathfrak{a}}_{x}\widehat{p}_{x}\widehat{\mathfrak{a}}_{x}\widehat{q}_{x}+\widehat{\mathfrak{a}}_{x}\widehat{p}_{x}\widehat{\mathfrak{a}}_{y}\widehat{q}_{y}+\widehat{\mathfrak{a}}_{x}\widehat{p}_{x}\widehat{\mathfrak{a}}_{z}\widehat{q}_{z}\\
 & +\widehat{\mathfrak{a}}_{y}\widehat{p}_{y}\widehat{\mathfrak{a}}_{x}\widehat{q}_{x}+\widehat{\mathfrak{a}}_{y}\widehat{p}_{y}\widehat{\mathfrak{a}}_{y}\widehat{q}_{y}+\widehat{\mathfrak{a}}_{y}\widehat{p}_{y}\widehat{\mathfrak{a}}_{z}\widehat{q}_{z}\\
 & +\widehat{\mathfrak{a}}_{z}\widehat{p}_{z}\widehat{\mathfrak{a}}_{x}\widehat{q}_{x}+\widehat{\mathfrak{a}}_{z}\widehat{p}_{z}\widehat{\mathfrak{a}}_{y}\widehat{q}_{y}+\widehat{\mathfrak{a}}_{z}\widehat{p}_{z}\widehat{\mathfrak{a}}_{z}\widehat{q}_{z}\\
= & \widehat{\mathfrak{a}}_{x}\widehat{\mathfrak{a}}_{x}\widehat{p}_{x}\widehat{q}_{x}+\widehat{\mathfrak{a}}_{x}\widehat{\mathfrak{a}}_{y}\widehat{p}_{x}\widehat{q}_{y}+\widehat{\mathfrak{a}}_{x}\widehat{\mathfrak{a}}_{z}\widehat{p}_{x}\widehat{q}_{z}\\
 & +\widehat{\mathfrak{a}}_{y}\widehat{\mathfrak{a}}_{x}\widehat{p}_{y}\widehat{q}_{x}+\widehat{\mathfrak{a}}_{y}\widehat{\mathfrak{a}}_{y}\widehat{p}_{y}\widehat{q}_{y}+\widehat{\mathfrak{a}}_{y}\widehat{\mathfrak{a}}_{z}\widehat{p}_{y}\widehat{q}_{z}\\
 & +\widehat{\mathfrak{a}}_{z}\widehat{\mathfrak{a}}_{x}\widehat{p}_{z}\widehat{q}_{x}+\widehat{\mathfrak{a}}_{z}\widehat{\mathfrak{a}}_{y}\widehat{p}_{z}\widehat{q}_{y}+\widehat{\mathfrak{a}}_{z}\widehat{\mathfrak{a}}_{z}\widehat{p}_{z}\widehat{q}_{z}
\end{align*}
}{\footnotesize\par}

{\footnotesize{}$\widehat{\boldsymbol{\mathfrak{a}}}$} is Second
Relativized Dirac Matrice(SRDM), {\footnotesize{}as for $\widehat{\mathfrak{a}}_{\mathrm{i}}=\mathrm{f_{i}}(\epsilon)\widehat{\alpha}_{\mathrm{i}}$,
it's easy to get the $\epsilon_{\mathrm{i}}$ as coefficeint of $\widehat{\sigma}_{\mathrm{i}}$
with the anti-commutate relation of Pauli matrices if $\mathrm{f_{i}}(\epsilon)$
satisfies ternary equations
\[
\mathrm{\forall\mathrm{i,j,k}:f_{i}}(\epsilon)\mathrm{f_{j}}(\epsilon)\equiv\epsilon_{\mathrm{k}}
\]
it gives the only result $\widehat{\mathfrak{a}}_{\mathrm{i}}=\pm\sqrt{\frac{\epsilon_{\mathrm{j}}\epsilon_{\mathrm{k}}}{\epsilon_{\mathrm{i}}}}\widehat{\alpha}_{\mathrm{i}}$,
positive one will be adopted.}{\footnotesize\par}
\end{theorem}

\end{frame}
\pagebreak{}
\begin{frame}{1.2 \textbf{Back to SRTP-fpFW Hamiltonian}}

$\mathrm{f_{i}}(\epsilon)$ in previous derivation can be generalized
to non-identity matrix operator, since the matrice commutate with
all alpha matrices has the form: 
\[
\mathrm{f_{i}}(\epsilon)=\left(\begin{array}{cccc}
\widehat{m}_{\mathrm{i}} &  & \widehat{n}_{\mathrm{i}}\\
 & \widehat{m}_{\mathrm{i}} &  & \widehat{n}_{\mathrm{i}}\\
\widehat{n}_{\mathrm{i}} &  & \widehat{m}_{\mathrm{i}}\\
 & \widehat{n}_{\mathrm{i}} &  & \widehat{m}_{\mathrm{i}}
\end{array}\right)
\]
which $m$ and $n$ are arbitrary operators. The reason why we have
to generalize $\mathrm{f_{i}}(\epsilon)$ is that current $\widehat{\mathfrak{a}}_{\mathrm{i}}$
do not satisfies the identity property of Dirac matrices, this makes
the scalar term in Pauli vector rule changed from original Darwin
term, which is not expected in SRTP Hamiltonian (SRTP is to modify
the spin-dependent Thomas precession). Furthermore, unidentity $\widehat{\mathfrak{a}}_{\mathrm{i}}$
makes the Dirac equation consists of $\widehat{\mathfrak{a}}_{\mathrm{i}}$
no longer satisfies Lorentz invarient since the energy-momentum relationship
are destroyed. Construct two sets of equations: 
\[
\forall\mathrm{i}:\begin{cases}
\left\{ \widehat{m}_{\mathrm{i}},\widehat{n}_{\mathrm{i}}\right\} =0\\
\widehat{m}_{\mathrm{i}}^{2}+\widehat{n}_{\mathrm{i}}^{2}=1
\end{cases}
\]

\[
\forall\mathrm{i,j,k}:\begin{cases}
\widehat{m}_{\mathrm{i}}\widehat{n}_{\mathrm{j}}+\widehat{n}_{\mathrm{i}}\widehat{m}_{\mathrm{j}}=0\\
\widehat{m}_{\mathrm{i}}\widehat{m}_{\mathrm{j}}+\widehat{n}_{\mathrm{i}}\widehat{n}_{\mathrm{j}}=\epsilon_{k}
\end{cases}
\]
first equations give $(\widehat{m}_{\mathrm{i}}+\widehat{n}_{\mathrm{i}})^{2}=1$,
second equations give $\widehat{m}_{\mathrm{i}}+\widehat{n}_{\mathrm{i}}=\pm\sqrt{\frac{\epsilon_{\mathrm{j}}\epsilon_{\mathrm{k}}}{\epsilon_{\mathrm{i}}}}$,
they are contradictory, therefore we can't find a fully consistent
SRDM, the only way is a hybrid treatment of scalar term and spin-dependent
term, except, magnetic SOC term will follow scalar term since the
energy-momentum relationship won't change.
\end{frame}
\pagebreak{}
\begin{frame}{1.2 \textbf{Back to SRTP-fpFW Hamiltonian}}

Back to Eq. (11), (12), further 4-component coupling gives{\footnotesize{}
\begin{align*}
\widehat{\varepsilon}_{\mathrm{1}}^{1e}= & \widehat{\varepsilon}^{1e}-\frac{1}{4}\left[\frac{\epsilon+c^{2}}{\sqrt{\epsilon\left(\epsilon+c^{2}\right)}},\left[\frac{\epsilon+c^{2}}{\sqrt{\epsilon\left(\epsilon+c^{2}\right)}},\widehat{\varepsilon}^{1e}\right]\right]+\frac{1}{4}\left[\frac{\beta\mathcal{O}}{\sqrt{\epsilon\left(\epsilon+c^{2}\right)}},\left[\frac{\beta\mathcal{O}}{\sqrt{\epsilon\left(\epsilon+c^{2}\right)}},\widehat{\varepsilon}^{1e}\right]\right]\\
\widehat{\varepsilon}_{\mathrm{1}}^{2e}= & \widehat{\varepsilon}^{2e}-\frac{1}{4}\left[\frac{\epsilon+c^{2}}{\sqrt{\epsilon\left(\epsilon+c^{2}\right)}},\left[\frac{\epsilon+c^{2}}{\sqrt{\epsilon\left(\epsilon+c^{2}\right)}},\widehat{\varepsilon}^{2e}\right]\right]+\frac{1}{4}\left[\frac{\beta\mathcal{O}}{\sqrt{\epsilon\left(\epsilon+c^{2}\right)}},\left[\frac{\beta\mathcal{O}}{\sqrt{\epsilon\left(\epsilon+c^{2}\right)}},\widehat{\varepsilon}^{2e}\right]\right]
\end{align*}
use $ABA=\frac{1}{2}\left\{ A^{2},B\right\} -\left[A,\left[A,B\right]\right]$,
the odd term still satisfy $\beta\mathcal{O}=-\mathcal{O}\beta$ and
$\left[A_{p},\mathcal{O}\right]=0$, so SRTP-fpFW Hamiltonian has
still the traditional form:}{\footnotesize\par}

{\footnotesize{}
\begin{align}
\varepsilon_{0} & =\beta\epsilon\\
\varepsilon_{1} & =A_{p}\left(V+R_{p}VR_{p}\right)A_{p}\\
\mathcal{O}_{1} & =\beta A_{p}\left[R_{p},V\right]A_{p}\\
\varepsilon_{1}^{2e-C} & =\underset{\mathrm{i}<\mathrm{j}}{\sum}A_{\mathrm{i}}A_{\mathrm{j}}\left(V_{\mathrm{ij}}+R_{\mathrm{i}}V_{\mathrm{ij}}R_{\mathrm{i}}+R_{\mathrm{j}}V_{\mathrm{ij}}R_{\mathrm{j}}+R_{\mathrm{i}}R_{\mathrm{j}}V_{\mathrm{ij}}R_{\mathrm{i}}R_{\mathrm{j}}\right)A_{\mathrm{i}}A_{\mathrm{j}}\\
\varepsilon_{1}^{2e-B} & =\underset{\mathrm{i}<\mathrm{j}}{\sum}A_{\mathrm{i}}A_{\mathrm{j}}\left(R_{\mathrm{i}}R_{\mathrm{j}}B_{\mathrm{ij}}+R_{\mathrm{i}}B_{\mathrm{ij}}R_{\mathrm{j}}+R_{\mathrm{j}}B_{\mathrm{ij}}R_{\mathrm{i}}+B_{\mathrm{ij}}R_{\mathrm{i}}R_{\mathrm{j}}\right)A_{\mathrm{i}}A_{\mathrm{j}}
\end{align}
}{\footnotesize\par}

{\footnotesize{}which Eq. (13) is kinetic term, $\epsilon=\sqrt{c^{4}+(\hat{\mathbf{\sigma}}\cdot\widehat{\boldsymbol{\pi}})^{2}c^{2}}$;
Eq. (14) is first-order even term, $A_{p}=\sqrt{\frac{\epsilon+c^{2}}{2\epsilon}}$,
$R_{p}=\frac{c\mathfrak{\boldsymbol{a}}\cdot\boldsymbol{\pi}}{\epsilon+c^{2}}$;
Eq. (15) is first-order odd term, Eq. (16) is two-electron Coulomb
interaction, Eq. (17) is one-photon exchange correction (Breit term),
$B_{\mathrm{ij}}=-\frac{1}{2r_{\mathrm{ij}}}\left(\mathfrak{s}_{\mathrm{i}}\cdot\mathfrak{s}_{\mathrm{j}}+\left(\mathfrak{s}_{\mathrm{i}}\cdot\bar{r}_{\mathrm{ij}}\right)\left(\mathfrak{s}_{\mathrm{j}}\cdot\bar{r}_{\mathrm{ij}}\right)\right)$.}{\footnotesize\par}
\end{frame}
\pagebreak{}
\begin{frame}{1.2 \textbf{Back to SRTP-fpFW Hamiltonian}}

The lead term of SRTP correction is in $c^{-4}$ order, so it is necessary
to consider the radiative correction (in order of $c^{-3},c^{-4}$,
spin-dependent) phenomenologically, which is the anomalous magnetic
moments of electron. {\footnotesize{}According to Douglas and Kroll
(Ann. Phys., 1974, 82: 89-155), }radiative correction only works on
spin-dependent terms (not include scalar terms), so we apply the same
treatment to SRTP correction and radiative correction, this makes
the $\varepsilon_{1}$ term becomes

{\footnotesize{}
\begin{align}
\varepsilon_{1}^{R} & =A_{p}\left(V+R_{p}^{R}VR_{p}^{R}\right)A_{p}\\
\varepsilon_{1}^{2e-C} & =\underset{\mathrm{i}<\mathrm{j}}{\sum}A_{\mathrm{i}}A_{\mathrm{j}}\left(V_{\mathrm{ij}}+R_{\mathrm{i}}^{R}V_{\mathrm{ij}}R_{\mathrm{i}}^{R}+R_{\mathrm{j}}^{R}V_{\mathrm{ij}}R_{\mathrm{j}}^{R}+R_{\mathrm{i}}^{R}R_{\mathrm{j}}^{R}V_{\mathrm{ij}}R_{\mathrm{i}}^{R}R_{\mathrm{j}}^{R}\right)A_{\mathrm{i}}A_{\mathrm{j}}\\
\varepsilon_{1}^{2e-B} & =\underset{\mathrm{i}<\mathrm{j}}{\sum}A_{\mathrm{i}}A_{\mathrm{j}}\left(R_{\mathrm{i}}^{R}R_{\mathrm{j}}^{R}B_{\mathrm{ij}}^{R}+R_{\mathrm{i}}^{R}B_{\mathrm{ij}}^{R}R_{\mathrm{j}}^{R}+R_{\mathrm{j}}^{R}B_{\mathrm{ij}}^{R}R_{\mathrm{i}}^{R}+B_{\mathrm{ij}}^{R}R_{\mathrm{i}}^{R}R_{\mathrm{j}}^{R}\right)A_{\mathrm{i}}A_{\mathrm{j}}
\end{align}
}{\footnotesize\par}

{\footnotesize{}which $R_{p}^{R}=\frac{c\mathfrak{\boldsymbol{a}}^{R}\cdot\boldsymbol{\pi}}{\epsilon+c^{2}},\mathfrak{\boldsymbol{a}}_{\mathrm{i}}^{R}=\left(\sqrt{\frac{\epsilon_{\mathrm{j}}\epsilon_{\mathrm{k}}}{\epsilon_{\mathrm{i}}}}+\sqrt{\frac{1}{2\pi c}-0.328\frac{1}{\pi^{2}c^{2}}}\right)\mathfrak{\alpha}_{\mathrm{i}};B_{\mathrm{ij}}^{R}=-\frac{1}{2r_{\mathrm{ij}}}\left(\mathfrak{s}_{\mathrm{i}}^{R}\cdot\mathfrak{s}_{\mathrm{j}}^{R}+\left(\mathfrak{s}_{\mathrm{i}}^{R}\cdot\bar{r}_{\mathrm{ij}}\right)\left(\mathfrak{s}_{\mathrm{j}}^{R}\cdot\bar{r}_{\mathrm{ij}}\right)\right),\mathfrak{s}_{\mathrm{x}}^{R}=\left(\epsilon_{\mathrm{x}}+\frac{1}{2\pi c}-0.328\frac{1}{\pi^{2}c^{2}}\right)\sigma_{\mathrm{x}}$.
}The cross term between SRTP correction and anomalous magnetic moment
correction is not needed when higher order terms are not considered.
In addition, other QED effects such as multiphoton exchange also have
an effect on spin-dependent terms, which will not be considered in
this paper.
\end{frame}
\pagebreak{}
\begin{frame}{Contents}
\begin{description}
\item [{\textcolor{lightgray}{1}}] \textcolor{lightgray}{Derivation of
}\textbf{SRTP-fpFW}\textcolor{lightgray}{{} Hamiltonian}
\begin{description}
\item [{\textcolor{lightgray}{1.1}}] \textcolor{lightgray}{Operatorization
of SRTP Hamiltonian}
\item [{\textcolor{lightgray}{1.2}}] \textbf{Back to SRTP-fpFW Hamiltonian}
\end{description}
\item [{2}] \textbf{DKH transformation of SRTP-fpFW Hamiltonian}
\begin{description}
\item [{2.1}] \textbf{SRTP-DKH2-1e Hamiltonian}
\item [{2.2}] \textbf{SRTP-DKH2-2e Hamiltonian}
\end{description}
\item [{\textcolor{lightgray}{3}}] \textcolor{lightgray}{Modified DKH Hamiltonian}
\begin{description}
\item [{\textcolor{lightgray}{3.1}}] \textcolor{lightgray}{Operatorization
of SRTP Hamiltonian}
\item [{\textcolor{lightgray}{3.2}}] \textcolor{lightgray}{Operatorization
of SRTP Hamiltonian}
\end{description}
\begin{block}{Agreements}
\begin{itemize}
\item {\small{}All derication is based on a.u., the charge and mass of electron
is equal to 1, the magnitude of spin angular momentum is equal to
1/2. }{\small\par}
\item {\small{}Classical vector is represented as $\overrightarrow{\mathrm{a}}$;
vector operator is represented as $\mathbf{\hat{\mathbf{a}}}$.}{\small\par}
\end{itemize}
\end{block}
\end{description}
\end{frame}
\pagebreak{}
\begin{frame}{\textbf{SRTP-DKH2-1e Hamiltonian}}

For ease of the subsequent derivation, the one-electron SRTP-fpFW
Hamiltonian must be expressed in terms of momentum representation,
the kernel of the momentum representation of $\varepsilon_{1}$ and
$\mathcal{O}_{1}$ should be expressed as
\begin{align}
\varepsilon_{1}\left(\boldsymbol{p},\boldsymbol{p}^{\prime}\right) & =A_{p}V_{pp^{\prime}}A_{p^{\prime}}+A_{p}R_{p}V_{pp^{\prime}}R_{p^{\prime}}A_{p^{\prime}}\\
\mathcal{O}_{1}\left(\boldsymbol{p},\boldsymbol{p}^{\prime}\right) & =\beta A_{p}R_{p}V_{pp^{\prime}}A_{p^{\prime}}-\beta A_{p}V_{pp^{\prime}}R_{p^{\prime}}A_{p^{\prime}}\\
V_{pp^{\prime}} & =-\mathrm{Z}\sqrt{\frac{2}{\pi}}\frac{1}{\left(\boldsymbol{p}-\boldsymbol{p}^{\prime}\right)^{2}}\nonumber 
\end{align}

According to Hess(J. Chem. Phys., 2002, 117: 9215--9226)(Recent Advances
in Relativistic Molecular Theory, 2004, 137--190), the optimized
Douglas Kroll Hess transformation follows the form
\begin{align*}
H_{\mathrm{DKH2,1e}} & =\left(a_{0}+\stackrel[k=1]{\infty}{\sum}a_{k}W_{1}^{k}\right)\left(\varepsilon_{0}+\varepsilon_{1}+\mathcal{O}_{1}\right)\left(a_{0}+\stackrel[k=1]{\infty}{\sum}\left(-1\right)^{k}a_{k}W_{1}^{k}\right)\\
 & =\varepsilon_{0}+\varepsilon_{1}+\varepsilon_{2}+\varepsilon_{3}+\stackrel[\mathrm{k}=4]{\infty}{\sum}\mathcal{\epsilon}_{\mathrm{k}}^{(2)}+\stackrel[\mathrm{k}=4]{\infty}{\sum}\mathcal{O}_{\mathrm{k}}^{(2)}\\
W_{1} & =\frac{a_{0}}{a_{1}}\beta\frac{\mathcal{O}_{1}\left(\boldsymbol{p},\boldsymbol{p}^{\prime}\right)}{\epsilon+\epsilon^{\prime}}=\beta\frac{\mathcal{O}_{1}\left(\boldsymbol{p},\boldsymbol{p}^{\prime}\right)}{\epsilon+\epsilon^{\prime}}
\end{align*}
We use the optimal unitary transformation as DKH transformation, the
coefficient is set as $U^{\mathrm{opt}}$

\includegraphics{picture/DKcoefficient}
\end{frame}
\pagebreak{}
\begin{frame}{\textbf{SRTP-DKH2-1e Hamiltonian}}

After the DKH2 transformation, only $\varepsilon_{2}$ and $\varepsilon_{3}$
can be determined.\\
\begin{align*}
\varepsilon_{2}\varphi\left(\boldsymbol{p}\right)= & \frac{1}{2}\left[W_{1},\mathcal{O}_{1}\right]\varphi\left(\boldsymbol{p}\right)\\
= & \frac{1}{2}W_{1}\mathcal{O}_{1}\varphi\left(\boldsymbol{p}\right)-\frac{1}{2}\mathcal{O}_{1}W_{1}\varphi\left(\boldsymbol{p}\right)\\
= & \frac{1}{2}\left(\int\mathrm{d}\boldsymbol{p}^{\prime}\beta\frac{\mathcal{O}_{1}\left(\boldsymbol{p},\boldsymbol{p}^{\prime}\right)}{\epsilon+\epsilon^{\prime}}\int\mathrm{d}\boldsymbol{p}^{\prime\prime}\mathcal{O}_{1}\left(\boldsymbol{p}^{\prime},\boldsymbol{p}^{\prime\prime}\right)\varphi\left(\boldsymbol{p}^{\prime\prime}\right)\right)\\
 & -\frac{1}{2}\left(\int\mathrm{d}\boldsymbol{p}^{\prime}\mathcal{O}_{1}\left(\boldsymbol{p},\boldsymbol{p}^{\prime}\right)\int\mathrm{d}\boldsymbol{p}^{\prime\prime}\beta\frac{\mathcal{O}_{1}\left(\boldsymbol{p}^{\prime},\boldsymbol{p}^{\prime\prime}\right)}{\epsilon^{\prime}+\epsilon^{\prime\prime}}\varphi\left(\boldsymbol{p}^{\prime\prime}\right)\right)\\
= & \frac{1}{2}\left(\iint\mathrm{d}\boldsymbol{p}^{\prime}\mathrm{d}\boldsymbol{p}^{\prime\prime}\left[\beta\frac{\mathcal{O}_{1}\left(\boldsymbol{p},\boldsymbol{p}^{\prime}\right)}{\epsilon+\epsilon^{\prime}}\mathcal{O}_{1}\left(\boldsymbol{p}^{\prime},\boldsymbol{p}^{\prime\prime}\right)-\mathcal{O}_{1}\left(\boldsymbol{p},\boldsymbol{p}^{\prime}\right)\beta\frac{\mathcal{O}_{1}\left(\boldsymbol{p}^{\prime},\boldsymbol{p}^{\prime\prime}\right)}{\epsilon^{\prime}+\epsilon^{\prime\prime}}\right]\varphi\left(\boldsymbol{p}^{\prime\prime}\right)\right)
\end{align*}
its kernel reads
\begin{align*}
\varepsilon_{2}\left(\boldsymbol{p},\boldsymbol{p}^{\prime},\boldsymbol{p}^{\prime\prime}\right)= & \beta\frac{\mathcal{O}_{1}\left(\boldsymbol{p},\boldsymbol{p}^{\prime}\right)}{2\left(\epsilon+\epsilon^{\prime}\right)}\mathcal{O}_{1}\left(\boldsymbol{p}^{\prime},\boldsymbol{p}^{\prime\prime}\right)-\mathcal{O}_{1}\left(\boldsymbol{p},\boldsymbol{p}^{\prime}\right)\beta\frac{\mathcal{O}_{1}\left(\boldsymbol{p}^{\prime},\boldsymbol{p}^{\prime\prime}\right)}{2\left(\epsilon^{\prime}+\epsilon^{\prime\prime}\right)}\\
= & \frac{\left(A_{p}R_{p}V_{pp^{\prime}}A_{p^{\prime}}-A_{p}V_{pp^{\prime}}R_{p^{\prime}}A_{p^{\prime}}\right)\left(\beta A_{p^{\prime}}R_{p^{\prime}}V_{p^{\prime}p^{\prime\prime}}A_{p^{\prime\prime}}-\beta A_{p^{\prime}}V_{p^{\prime}p^{\prime\prime}}R_{p^{\prime\prime}}A_{p^{\prime\prime}}\right)}{2\left(\epsilon+\epsilon^{\prime}\right)}\\
 & -\frac{\left(\beta A_{p}R_{p}V_{pp^{\prime}}A_{p^{\prime}}-\beta AV_{pp^{\prime}}R_{p^{\prime}}A_{p^{\prime}}\right)\left(A_{p^{\prime}}R_{p^{\prime}}V_{p^{\prime}p^{\prime\prime}}A_{p^{\prime\prime}}-A_{p^{\prime}}V_{p^{\prime}p^{\prime\prime}}R_{p^{\prime\prime}}A_{p^{\prime\prime}}\right)}{2\left(\epsilon^{\prime}+\epsilon^{\prime\prime}\right)}
\end{align*}
For the case of no external magnetic field (excluding Breit term),
$\left[A_{p},R_{p}\right]=0$, it gives
\begin{align*}
\varepsilon_{2}\left(\boldsymbol{p},\boldsymbol{p}^{\prime},\boldsymbol{p}^{\prime\prime}\right)= & -\beta\frac{A_{p}R_{p}V_{pp^{\prime}}R_{p^{\prime}}A_{p^{\prime}}A_{p^{\prime}}V_{p^{\prime}p^{\prime\prime}}A_{p^{\prime\prime}}}{2\left(\epsilon+\epsilon^{\prime}\right)}+\beta\frac{A_{p}R_{p}V_{pp^{\prime}}A_{p^{\prime}}A_{p^{\prime}}V_{p^{\prime}p^{\prime\prime}}R_{p^{\prime\prime}}A_{p^{\prime\prime}}}{2\left(\epsilon+\epsilon^{\prime}\right)}\\
 & +\beta\frac{A_{p}V_{pp^{\prime}}R_{p^{\prime}}A_{p^{\prime}}A_{p^{\prime}}R_{p^{\prime}}V_{p^{\prime}p^{\prime\prime}}A_{p^{\prime\prime}}}{2\left(\epsilon+\epsilon^{\prime}\right)}-\beta\frac{A_{p}V_{pp^{\prime}}A_{p^{\prime}}A_{p^{\prime}}R_{p^{\prime}}V_{p^{\prime}p^{\prime\prime}}R_{p^{\prime\prime}}A_{p^{\prime\prime}}}{2\left(\epsilon+\epsilon^{\prime}\right)}\\
 & -\beta\frac{A_{p}R_{p}V_{pp^{\prime}}R_{p^{\prime}}A_{p^{\prime}}A_{p^{\prime}}V_{p^{\prime}p^{\prime\prime}}A_{p^{\prime\prime}}}{2\left(\epsilon^{\prime}+\epsilon^{\prime\prime}\right)}+\beta\frac{A_{p}R_{p}V_{pp^{\prime}}A_{p^{\prime}}A_{p^{\prime}}V_{p^{\prime}p^{\prime\prime}}R_{p^{\prime\prime}}A_{p^{\prime\prime}}}{2\left(\epsilon^{\prime}+\epsilon^{\prime\prime}\right)}\\
 & +\beta\frac{A_{p}V_{pp^{\prime}}R_{p^{\prime}}A_{p^{\prime}}A_{p^{\prime}}R_{p^{\prime}}V_{p^{\prime}p^{\prime\prime}}A_{p^{\prime\prime}}}{2\left(\epsilon^{\prime}+\epsilon^{\prime\prime}\right)}-\beta\frac{A_{p}V_{pp^{\prime}}A_{p^{\prime}}A_{p^{\prime}}R_{p^{\prime}}V_{p^{\prime}p^{\prime\prime}}R_{p^{\prime\prime}}A_{p^{\prime\prime}}}{2\left(\epsilon^{\prime}+\epsilon^{\prime\prime}\right)}
\end{align*}
\end{frame}
\pagebreak{}
\begin{frame}{\textbf{SRTP-DKH2-1e Hamiltonian}}

\begin{align*}
\varepsilon_{3}\varphi\left(\boldsymbol{p}\right)= & \frac{1}{2}\left[W_{1},\left[W_{1},\varepsilon_{1}\right]\right]\varphi\left(\boldsymbol{p}\right)\\
= & \frac{1}{2}W_{1}W_{1}\varepsilon_{1}\varphi\left(\boldsymbol{p}\right)-W_{1}\varepsilon_{1}W_{1}\varphi\left(\boldsymbol{p}\right)+\frac{1}{2}\varepsilon_{1}W_{1}W_{1}\\
= & \frac{1}{2}\left(\int\mathrm{d}\boldsymbol{p}^{\prime}\beta\frac{\mathcal{O}_{1}\left(\boldsymbol{p},\boldsymbol{p}^{\prime}\right)}{\epsilon+\epsilon^{\prime}}\int\mathrm{d}\boldsymbol{p}^{\prime\prime}\beta\frac{\mathcal{O}_{1}\left(\boldsymbol{p}^{\prime},\boldsymbol{p}^{\prime\prime}\right)}{\epsilon^{\prime}+\epsilon^{\prime\prime}}\int\mathrm{d}\boldsymbol{p}^{\prime\prime\prime}\varepsilon_{1}\left(\boldsymbol{p}^{\prime\prime},\boldsymbol{p}^{\prime\prime\prime}\right)\varphi\left(\boldsymbol{p}^{\prime\prime\prime}\right)\right)\\
 & -\left(\int\mathrm{d}\boldsymbol{p}^{\prime}\beta\frac{\mathcal{O}_{1}\left(\boldsymbol{p},\boldsymbol{p}^{\prime}\right)}{\epsilon+\epsilon^{\prime}}\int\mathrm{d}\boldsymbol{p}^{\prime\prime}\varepsilon_{1}\left(\boldsymbol{p}^{\prime},\boldsymbol{p}^{\prime\prime}\right)\int\mathrm{d}\boldsymbol{p}^{\prime\prime\prime}\beta\frac{\mathcal{O}_{1}\left(\boldsymbol{p}^{\prime\prime},\boldsymbol{p}^{\prime\prime\prime}\right)}{\epsilon^{\prime\prime}+\epsilon^{\prime\prime\prime}}\varphi\left(\boldsymbol{p}^{\prime\prime\prime}\right)\right)\\
 & +\frac{1}{2}\left(\int\mathrm{d}\boldsymbol{p}^{\prime}\varepsilon_{1}\left(\boldsymbol{p},\boldsymbol{p}^{\prime}\right)\int\mathrm{d}\boldsymbol{p}^{\prime\prime}\beta\frac{\mathcal{O}_{1}\left(\boldsymbol{p}^{\prime},\boldsymbol{p}^{\prime\prime}\right)}{\epsilon^{\prime}+\epsilon^{\prime\prime}}\int\mathrm{d}\boldsymbol{p}^{\prime\prime\prime}\beta\frac{\mathcal{O}_{1}\left(\boldsymbol{p}^{\prime\prime},\boldsymbol{p}^{\prime\prime\prime}\right)}{\epsilon^{\prime\prime}+\epsilon^{\prime\prime\prime}}\varphi\left(\boldsymbol{p}^{\prime\prime\prime}\right)\right)
\end{align*}
its kernel reads
\begin{align*}
\varepsilon_{3}\left(\boldsymbol{p},\boldsymbol{p}^{\prime},\boldsymbol{p}^{\prime\prime},\boldsymbol{p}^{\prime\prime\prime}\right)= & \beta\frac{\mathcal{O}_{1}\left(\boldsymbol{p},\boldsymbol{p}^{\prime}\right)}{2\left(\epsilon+\epsilon^{\prime}\right)}\beta\frac{\mathcal{O}_{1}\left(\boldsymbol{p}^{\prime},\boldsymbol{p}^{\prime\prime}\right)}{\epsilon^{\prime}+\epsilon^{\prime\prime}}\varepsilon_{1}\left(\boldsymbol{p}^{\prime\prime},\boldsymbol{p}^{\prime\prime\prime}\right)\\
 & -\beta\frac{\mathcal{O}_{1}\left(\boldsymbol{p},\boldsymbol{p}^{\prime}\right)}{\epsilon+\epsilon^{\prime}}\varepsilon_{1}\left(\boldsymbol{p}^{\prime},\boldsymbol{p}^{\prime\prime}\right)\beta\frac{\mathcal{O}_{1}\left(\boldsymbol{p}^{\prime\prime},\boldsymbol{p}^{\prime\prime\prime}\right)}{\epsilon^{\prime\prime}+\epsilon^{\prime\prime\prime}}\\
 & +\varepsilon_{1}\left(\boldsymbol{p},\boldsymbol{p}^{\prime}\right)\beta\frac{\mathcal{O}_{1}\left(\boldsymbol{p}^{\prime},\boldsymbol{p}^{\prime\prime}\right)}{2\left(\epsilon^{\prime}+\epsilon^{\prime\prime}\right)}\beta\frac{\mathcal{O}_{1}\left(\boldsymbol{p}^{\prime\prime},\boldsymbol{p}^{\prime\prime\prime}\right)}{\epsilon^{\prime\prime}+\epsilon^{\prime\prime\prime}}
\end{align*}
Clearly $\varepsilon_{2}$ is of order $c^{-4}$ and $\varepsilon_{3}$
is of order $c^{-6}$, we have no need to consider terms higher than
order $c^{-4}$. Therefore, for one-electron term, only $\varepsilon_{0},\varepsilon_{1},\varepsilon_{2}$
is considered, and $\varepsilon_{2}$ does not take into account the
influence of SRTP effect.
\end{frame}
\pagebreak{}
\begin{frame}{\textbf{DKH2-2e Hamiltonian}}

For the DKHn Hamiltonian, the two-electron DKH transformation entail
much smaller effects on the relative energy and electronic structure
than the one-electron DKH transformation, despit they have the same
c-order. Considering the complexity of the two-electron integrals,
we push the two-electron term only to order $c^{-2}$, which means
that only the fpFW-2e (Coulomb and Breit) term and the DKH2-2e (Coulomb)
term are considered.
\begin{align}
\varepsilon_{\mathrm{1,2e},C} & =\underset{i<j}{\sum}A_{i}A_{j}\left[\frac{1}{r_{ij}}+R_{i}\frac{1}{r_{ij}}R_{i}+R_{j}\frac{1}{r_{ij}}R_{j}\right]A_{i}A_{j}\\
\varepsilon_{\mathrm{1,2e},B} & =\underset{i<j}{\sum}A_{i}A_{j}\left[R_{i}R_{j}\widetilde{B}_{ij}+R_{i}\widetilde{B}_{ij}R_{j}+R_{j}\widetilde{B}_{ij}R_{i}+\widetilde{B}_{ij}R_{i}R_{j}\right]A_{i}A_{j}
\end{align}

which $\widetilde{B}_{ij}=-\frac{1}{2r_{ij}}\left(\boldsymbol{\sigma}_{i}\cdot\boldsymbol{\sigma}_{j}+\left(\boldsymbol{\sigma}_{i}\cdot\hat{\boldsymbol{r}}_{ij}\right)\left(\boldsymbol{\sigma}_{j}\cdot\hat{\boldsymbol{r}}_{ij}\right)\right)$,{\footnotesize{}$A=\sqrt{\frac{\epsilon+c^{2}}{2\epsilon}}$,
$R=\frac{c\mathfrak{\boldsymbol{\alpha}}\cdot\boldsymbol{p}}{\epsilon+c^{2}}$,
this is $c^{-2}$ order truncation of Eq. (19) and Eq. (20). Eq. (23)
lacks the cross terms between }\textrm{$R_{i}$}{\footnotesize{} and
}\textrm{$R_{i}$}{\footnotesize{} because Dirac's equation indicates
that spin-same-orbital coupling, as the only $c^{-2}$ term in the
spin-dependent part of the 2e Coulomb interaction, occurs only in
non-cross terms. It is noteworthy that for 4-indicator integrals,
both Eq. (23) and Eq. (24) operate on the so-called \textquoteleft physical
representation\textquoteright{} rather than the \textquoteleft chemical
representation\textquoteright , for example}{\footnotesize\par}

\[
\left\langle p_{i}V_{ij}p_{j}\right\rangle _{\mu\nu\kappa\lambda}=\left\langle \chi_{\mu}\left(i\right)\chi_{\nu}\left(j\right)\right|\hat{p}_{i}\hat{V}_{ij}\hat{p}_{j}\left|\chi_{\kappa}\left(i\right)\chi_{\lambda}\left(j\right)\right\rangle 
\]

which shows the operators on both sides operate only on one of the
two Gaussian basis functions on each side, directly yielding{\footnotesize{}
4-indicator integrals} that can be converted into classical and exchange
terms of {\footnotesize{}2-indicator integrals}.
\end{frame}
\pagebreak{}
\begin{frame}{\textbf{DKH2-2e Hamiltonian}}

{\footnotesize{}It is noteworthy that for 4-indicator integrals, both
Eq. (23) and Eq. (24) operate on the so-called \textquoteleft physical
representation\textquoteright{} rather than the \textquoteleft chemical
representation\textquoteright , for example}{\footnotesize\par}

\[
\left\langle p_{i}V_{ij}p_{i}\right\rangle _{\mu\nu\kappa\lambda}=\left\langle \chi_{\mu}\left(i\right)\chi_{\nu}\left(j\right)\right|\hat{p}_{i}\hat{V}_{ij}\hat{p}_{i}\left|\chi_{\kappa}\left(i\right)\chi_{\lambda}\left(j\right)\right\rangle 
\]

which shows the operators on both sides operate only on one of the
two Gaussian basis functions on each side, directly yielding{\footnotesize{}
4-indicator integrals} that can be converted into classical and exchange
terms of {\footnotesize{}2-indicator integrals}.

Additionally, attention should be paid to how Hess's RI technique
can be applied to the 2-electron DKH2 integrals. In this case, we
need to use $p^{2}$ eigenbasis $\left|\chi^{p2}\right\rangle $,
insert the 2-electron identity

\begin{align*}
\left\langle A_{i}A_{j}R_{i}\frac{1}{r_{ij}}R_{i}A_{i}A_{j}\right\rangle _{\mu\nu\kappa\lambda}^{p2}= & \left\langle \chi_{\mu}^{p2}\left(i\right)\chi_{\nu}^{p2}\left(j\right)\right|A_{i}A_{j}R_{i}\frac{1}{r_{ij}}R_{i}A_{i}A_{j}\left|\chi_{\kappa}^{p2}\left(i\right)\chi_{\lambda}^{p2}\left(j\right)\right\rangle \\
= & \underset{mn}{\sum}\underset{pq}{\sum}\left\langle \chi_{\mu}^{p2}\left(i\right)\chi_{\nu}^{p2}\left(j\right)\right|A_{i}A_{j}\left|\chi_{m}^{p2}\left(i\right)\chi_{n}^{p2}\left(j\right)\right\rangle \left\langle \chi_{m}^{p2}\left(i\right)\chi_{n}^{p2}\left(j\right)\right|R_{i}\frac{1}{r_{ij}}R_{i}\left|\chi_{p}^{p2}\left(i\right)\chi_{q}^{p2}\left(j\right)\right\rangle \left\langle \chi_{p}^{p2}\left(i\right)\chi_{q}^{p2}\left(j\right)\right|A_{i}A_{j}\left|\chi_{\kappa}^{p2}\left(i\right)\chi_{\lambda}^{p2}\left(j\right)\right\rangle \\
= & \underset{mn}{\sum}\underset{pq}{\sum}\left\langle \chi_{\mu}^{p2}\left(i\right)\right|A_{i}\left|\chi_{m}^{p2}\left(i\right)\right\rangle \left\langle \chi_{\nu}^{p2}\left(j\right)\right|A_{j}\left|\chi_{n}^{p2}\left(j\right)\right\rangle \left\langle \chi_{m}^{p2}\left(i\right)\chi_{n}^{p2}\left(j\right)\right|R_{i}\frac{1}{r_{ij}}R_{i}\left|\chi_{p}^{p2}\left(i\right)\chi_{q}^{p2}\left(j\right)\right\rangle \left\langle \chi_{q}^{p2}\left(j\right)\right|A_{j}\left|\chi_{\lambda}^{p2}\left(j\right)\right\rangle \left\langle \chi_{p}^{p2}\left(i\right)\right|A_{i}\left|\chi_{\kappa}^{p2}\left(i\right)\right\rangle \\
= & A_{\mu}^{p2}\left(i\right)A_{\nu}^{p2}\left(j\right)\left\langle \chi_{\mu}^{p2}\left(i\right)\chi_{\nu}^{p2}\left(j\right)\right|R_{i}\frac{1}{r_{ij}}R_{i}\left|\chi_{\kappa}^{p2}\left(i\right)\chi_{\lambda}^{p2}\left(j\right)\right\rangle A_{\kappa}^{p2}\left(i\right)A_{\lambda}^{p2}\left(j\right)
\end{align*}

which $\left\langle \chi_{\mu}^{p2}\left(i\right)\right|A_{i}\left|\chi_{m}^{p2}\left(i\right)\right\rangle =A_{\mu}^{p2}\left(i\right)\delta_{\mu m}$,
the {\footnotesize{}4-indicator integral} of $p^{2}$ eigenbasis in
this expression is difficult to handle, we need to investigate their
contribution to the {\footnotesize{}2-indicator} Coulomb integral

\begin{align*}
\left\langle A_{i}A_{j}R_{i}\frac{1}{r_{ij}}R_{i}A_{i}A_{j}\right\rangle _{\mu\kappa,C}^{p2} & =\underset{\nu\lambda}{\sum}\left[\left\langle \chi_{\mu}^{p2}\left(i\right)\chi_{\nu}^{p2}\left(j\right)\right|A_{i}A_{j}R_{i}\frac{1}{r_{ij}}R_{i}A_{i}A_{j}\left|\chi_{\kappa}^{p2}\left(i\right)\chi_{\lambda}^{p2}\left(j\right)\right\rangle D_{\nu\lambda}^{p2}\right]\\
 & =A_{\mu}^{p2}\left(i\right)A_{\kappa}^{p2}\left(i\right)\underset{\nu\lambda}{\sum}\left[\left\langle \chi_{\mu}^{p2}\left(i\right)\chi_{\nu}^{p2}\left(j\right)\right|R_{i}\frac{1}{r_{ij}}R_{i}\left|\chi_{\kappa}^{p2}\left(i\right)\chi_{\lambda}^{p2}\left(j\right)\right\rangle D_{\nu\lambda}^{p2}A_{\nu}^{p2}\left(j\right)A_{\lambda}^{p2}\left(j\right)\right]
\end{align*}
\end{frame}
\pagebreak{}
\begin{frame}{\textbf{DKH2-2e Hamiltonian}}

which $D^{p2}$ is the single-electron density matrix in $p^{2}$
eigenbasis, let $D_{\nu\lambda}^{p2\prime}=D_{\nu\lambda}^{p2}A_{\nu}^{p2}\left(j\right)A_{\lambda}^{p2}\left(j\right)$,
for 2-electron integral $X_{ij}$, we have $\left\langle X_{ij}\right\rangle _{\mu\kappa,C}^{p2}=\left[\Omega^{\dagger}\left\langle X_{ij}\right\rangle _{C}\Omega\right]_{\mu\kappa}$,
so we get

\begin{align*}
\left\langle A_{i}A_{j}R_{i}\frac{1}{r_{ij}}R_{i}A_{i}A_{j}\right\rangle _{\mu\kappa,C}^{p2} & =A_{\mu}^{p2}\left(i\right)A_{\kappa}^{p2}\left(i\right)\underset{\nu\lambda}{\sum}\left[\left\langle \chi_{\mu}^{p2}\left(i\right)\chi_{\nu}^{p2}\left(j\right)\right|R_{i}\frac{1}{r_{ij}}R_{i}\left|\chi_{\kappa}^{p2}\left(i\right)\chi_{\lambda}^{p2}\left(j\right)\right\rangle D_{\nu\lambda}^{p2\prime}\right]\\
 & =A_{\mu}^{p2}\left(i\right)A_{\kappa}^{p2}\left(i\right)\left\langle R_{i}\frac{1}{r_{ij}}R_{i}\right\rangle _{\mu\kappa,C}^{p2\prime}\\
 & =A_{\mu}^{p2}\left(i\right)A_{\kappa}^{p2}\left(i\right)\left[\Omega^{\dagger}\left\langle R_{i}\frac{1}{r_{ij}}R_{i}\right\rangle _{C}^{\prime}\Omega\right]_{\mu\kappa}
\end{align*}

similarly, {\footnotesize{}for 2-indicator} exchange integral we get

\begin{align*}
\left\langle A_{i}A_{j}R_{i}\frac{1}{r_{ij}}R_{i}A_{i}A_{j}\right\rangle _{\mu\lambda,X}^{p2} & =A_{\mu}^{p2}\left(i\right)A_{\lambda}^{p2}\left(i\right)\underset{\nu\kappa}{\sum}\left[\left\langle \chi_{\mu}^{p2}\left(i\right)\chi_{\nu}^{p2}\left(j\right)\right|R_{i}\frac{1}{r_{ij}}R_{i}\left|\chi_{\lambda}^{p2}\left(i\right)\chi_{\kappa}^{p2}\left(j\right)\right\rangle D_{\nu\kappa}^{p2\prime}\right]\\
 & =A_{\mu}^{p2}\left(i\right)A_{\lambda}^{p2}\left(i\right)\left\langle R_{i}\frac{1}{r_{ij}}R_{i}\right\rangle _{\mu\lambda,X}^{p2\prime}\\
 & =A_{\mu}^{p2}\left(i\right)A_{\lambda}^{p2}\left(i\right)\left[\Omega^{\dagger}\left\langle R_{i}\frac{1}{r_{ij}}R_{i}\right\rangle _{X}^{\prime}\Omega\right]_{\mu\lambda}
\end{align*}
\end{frame}
\pagebreak{}
\begin{frame}{\textbf{DKH2-2e Hamiltonian}}

In order to truncate at $c^{-2}$, the two-electron DKH2 transformation
is performed only for the $\frac{1}{r_{ij}}$ term, then the first
order odd term is truncated to: 
\[
\mathcal{O}_{1,2e}=\underset{i<j}{\sum}A_{i}A_{j}\left(-\left[\frac{1}{r_{ij}},\beta_{i}R_{i}\right]-\left[\frac{1}{r_{ij}},\beta_{j}R_{j}\right]\right)A_{i}A_{j}
\]

thus

\begin{align}
H_{\mathrm{DKH2,2e}}= & \left(1+W_{1-2e}\right)\left(\varepsilon_{1-2e}+\mathcal{O}_{1-2e}\right)\left(1-W_{1-2e}\right)\\
= & \varepsilon_{1-2e}-\varepsilon_{1-2e}W_{1-2e}+\mathcal{O}_{1-2e}-\mathcal{O}_{1-2e}W_{1-2e}\nonumber \\
 & +W_{1-2e}\varepsilon_{1-2e}-W_{1-2e}\varepsilon_{1-2e}W_{1-2e}\nonumber \\
 & +W_{1-2e}\mathcal{O}_{1-2e}-W_{1-2e}\mathcal{O}_{1-2e}W_{1-2e}\nonumber 
\end{align}

$W_{1-2e}$ satisfies
\[
\mathcal{O}_{1,2e}+W_{1-2e}\varepsilon_{1-2e}-\varepsilon_{1-2e}W_{1-2e}=0
\]

$W_{1-2e}$ is given by
\[
W_{1,2e}=\beta_{i}\frac{\mathcal{O}_{1-2e}}{2g_{ij}}
\]

its even term reads(truncate at $c^{-2}$)
\begin{align}
\varepsilon_{\mathrm{2,2e}}= & W_{1-2e}\mathcal{O}_{1-2e}-\mathcal{O}_{1-2e}W_{1-2e}\nonumber \\
= & \frac{\mathcal{O}_{1-2e}^{2}}{g_{ij}}\nonumber \\
= & \beta_{i}\frac{A_{i}A_{j}\left(\left[\frac{1}{r_{ij}},\beta_{i}R_{i}\right]+\left[\frac{1}{r_{ij}},\beta_{j}R_{j}\right]\right)A_{i}A_{j}A_{i}A_{j}\left(\left[\frac{1}{r_{ij}},\beta_{i}R_{i}\right]+\left[\frac{1}{r_{ij}},\beta_{j}R_{j}\right]\right)A_{i}A_{j}}{g_{ij}}\nonumber \\
= & \beta_{i}\frac{A_{i}A_{j}\left(\frac{1}{r_{ij}}\beta_{i}R_{i}-\beta_{i}R_{i}\frac{1}{r_{ij}}+\frac{1}{r_{ij}}\beta_{j}R_{j}-\beta_{j}R_{j}\frac{1}{r_{ij}}\right)A_{i}A_{j}A_{i}A_{j}\left(\frac{1}{r_{ij}}\beta_{i}R_{i}-\beta_{i}R_{i}\frac{1}{r_{ij}}+\frac{1}{r_{ij}}\beta_{j}R_{j}-\beta_{j}R_{j}\frac{1}{r_{ij}}\right)A_{i}A_{j}}{g_{ij}}
\end{align}
To summarize, the two-electron DKH integral is

\[
H_{\mathrm{DKH2,2e}}=\varepsilon_{\mathrm{1-2e},C}+\varepsilon_{\mathrm{1-2e},B}+\varepsilon_{\mathrm{2-2e}}
\]
\end{frame}
\pagebreak
\end{document}
